% Part: normal-modal-logic
% Chapter: syntax-and-semantics
% Section: substitution

\documentclass[../../../include/open-logic-section]{subfiles}

\begin{document}

\olfileid[fr]{nml}{syn}{sub}

\olsection{Substitution simultanée}

Une \emph{instance} d'!!a{formula}~$!A$ s'obtient en remplaçant
toutes les occurrences d'!!a{propositional variable} dans~$!A$
par une autre !!{formula}. Nous parlerons souvent d'instances de
!!{formula}s, aussi bien à propos de la validité que de
l'!!{derivability}. Il est donc utile de définir précisément cette
notion.

\begin{defn}\ollabel{def:subst-inst}
  Soit $!A$ une !!{formula} modale. Soient $p_1$, \dots, $p_n$ des
  !!{propositional variable}s deux à deux distinctes et $!D_1$, \dots,
  $!D_n$ des !!{formula}s modales. Nous définissons
  $\SSubst{!A}{\subst{!D_1}{p_1}, \dots, \subst{!D_n}{p_n}}$ comme le
  résultat de la substitution simultanée de chaque $!D_i$ à~$p_i$
  dans~$!A$. Formellement, la définition se fait par induction
  sur~$!A$ :
  \begin{enumerate}
    \tagitem{prvFalse}{\indcase{!A}{\lfalse}{%
        $\SSubst{\indfrm}{\subst{!D_1}{p_1}, \dots,
          \subst{!D_n}{p_n}}$ vaut $\lfalse$.}}{}
    \tagitem{prvTrue}{\indcase{!A}{\ltrue}{%
        $\SSubst{\indfrm}{\subst{!D_1}{p_1}, \dots,
          \subst{!D_n}{p_n}}$ vaut $\ltrue$.}}{}
    \item \indcase{!A}{q}{$\SSubst{\indfrm}{\subst{!D_1}{p_1}, \dots,
        \subst{!D_n}{p_n}}$ vaut $q$, à condition que $q \not\ident p_i$ pour $i
      = 1$, \dots,~$n$.}
    \item \indcase{!A}{p_i}{$\SSubst{\indfrm}{\subst{!D_1}{p_1},
        \dots, \subst{!D_n}{p_n}}$ vaut $!D_i$.}
    \tagitem{prvNot}{\indcase{!A}{\lnot !B}{%
        $\SSubst{\indfrm}{\subst{!D_1}{p_1}, \dots, \subst{!D_n}{p_n}}$ vaut
        $\lnot \SSubst{!B}{\subst{!D_1}{p_1}, \dots, \subst{!D_n}{p_n}}$.}}{}
    \tagitem{prvAnd}{\indcase{!A}{(!B \land
        !C)}{$\SSubst{\indfrm}{\subst{!D_1}{p_1}, \dots,
          \subst{!D_n}{p_n}}$ vaut
        \[(\SSubst{!B}{\subst{!D_1}{p_1}, \dots, \subst{!D_n}{p_n}} \land
          \SSubst{!C}{\subst{!D_1}{p_1}, \dots, \subst{!D_n}{p_n}}).\]}}{}
    \tagitem{prvOr}{\indcase{!A}{(!B \lor
          !C)}{$\SSubst{\indfrm}{\subst{!D_1}{p_1}, \dots,
          \subst{!D_n}{p_n}}$ vaut
        \[(\SSubst{!B}{\subst{!D_1}{p_1}, \dots, \subst{!D_n}{p_n}} \lor
          \SSubst{!C}{\subst{!D_1}{p_1}, \dots, \subst{!D_n}{p_n}}).\]}}{}
    \tagitem{prvIf}{\indcase{!A}{(!B \lif
          !C)}{$\SSubst{\indfrm}{\subst{!D_1}{p_1}, \dots,
          \subst{!D_n}{p_n}}$ vaut
        \[(\SSubst{!B}{\subst{!D_1}{p_1}, \dots, \subst{!D_n}{p_n}} \lif
          \SSubst{!C}{\subst{!D_1}{p_1}, \dots, \subst{!D_n}{p_n}}).\]}}{}
    \tagitem{prvIff}{\indcase{!A}{(!B \liff
          !C)}{$\SSubst{\indfrm}{\subst{!D_1}{p_1}, \dots,
          \subst{!D_n}{p_n}}$ vaut
        \[(\SSubst{!B}{\subst{!D_1}{p_1}, \dots, \subst{!D_n}{p_n}} \liff
          \SSubst{!C}{\subst{!D_1}{p_1}, \dots, \subst{!D_n}{p_n}}).\]}}{}
    \tagitem{prvBox}{\indcase{!A}{\Box !B}{%
        $\SSubst{\indfrm}{\subst{!D_1}{p_1}, \dots, \subst{!D_n}{p_n}}$ vaut
        $\Box
        \SSubst{!B}{\subst{!D_1}{p_1}, \dots, \subst{!D_n}{p_n}}.$}}{}
    \tagitem{prvDiamond}{\indcase{!A}{\Diamond !B}{%
        $\SSubst{\indfrm}{\subst{!D_1}{p_1}, \dots, \subst{!D_n}{p_n}}$ vaut
        $\Diamond
        \SSubst{!B}{\subst{!D_1}{p_1}, \dots, \subst{!D_n}{p_n}}.$}}{}
  \end{enumerate}
  La !!{formula} $\SSubst{!A}{\subst{!D_1}{p_1}, \dots,
    \subst{!D_n}{p_n}}$ est appelée \emph{instance de substitution}
  de~$!A$.
\end{defn}

\begin{ex}
  Supposons que $!A$ soit $p_1 \lif \Box(p_1 \land p_2)$,
  que $!D_1$ soit $\Diamond(p_2 \lif p_3)$ et que $!D_2$
  soit $\lnot\Box p_1$. Alors
  $\SSubst{!A}{\subst{!D_1}{p_1}, \subst{!D_2}{p_2}}$ vaut
  \begin{align*}
    \Diamond(p_2 \lif p_3) &
    \lif \Box(\Diamond(p_2 \lif p_3) \land \lnot\Box p_1)
    \intertext{tandis que $\SSubst{!A}{\subst{!D_2}{p_1}, \subst{!D_1}{p_2}}$ vaut}
    \lnot\Box p_1 &
    \lif \Box(\lnot\Box p_1 \land \Diamond(p_2 \lif p_3))
    \intertext{La substitution simultanée diffère en général de la
      substitution successive.\newline Dans
      $\Subst{(\Subst{!A}{!D_1}{p_1})}{!D_2}{p_2}$, la seconde
      substitution modifie aussi les occurrences de $p_2$
      introduites par~$!D_1$; le résultat est :}
    \Diamond(\lnot\Box p_1 \lif p_3) &
    \lif \Box(\Diamond(\lnot\Box p_1
    \lif p_3) \land \lnot\Box p_1)\\
    \intertext{Avec
      $\Subst{(\Subst{!A}{!D_2}{p_2})}{!D_1}{p_1}$, la seconde
      substitution modifie aussi les occurrences de $p_1$
      introduites par~$!D_2$; nous obtenons :}
    \Diamond(p_2 \lif p_3) &
    \lif \Box(\Diamond(p_2
    \lif p_3) \land \lnot\Box\Diamond(p_2 \lif p_3)).
  \end{align*}
\end{ex}

\end{document}
